\documentclass[10pt,a4paper]{article}
\usepackage{amsmath,amssymb,amsthm,graphicx}
\usepackage{amsmath,amssymb,amsthm,hyperref,tikz,enumerate,subfigure}
\newtheorem{theorem}{Theorem}[section]
\newtheorem{theo}[theorem]{Theorem}
\newtheorem{lemm}[theorem]{Lemma}
\newtheorem{prop}[theorem]{Proposition}
\newtheorem{defi}[theorem]{Definition}
\newtheorem{rema}[theorem]{Remark}
\DeclareMathSymbol{\leqslant}{\mathalpha}{AMSa}{"36}
\DeclareMathSymbol{\geqslant}{\mathalpha}{AMSa}{"3E}
\DeclareMathSymbol{\eset}{\mathalpha}{AMSb}{"3F}
\renewcommand{\leq}{\;\leqslant\;}
\renewcommand{\geq}{\;\geqslant\;}
%\newcommand{\dd}{\,\text{\rm d}}




\newcommand{\la}{\label}
\newcommand{\Om}{\Omega}


\newcommand{\e}{\varepsilon}
\newcommand{\h}{\frac{1}{2}}

\newcommand{\rhoN}{\rho^{\pi}}
\newcommand{\si}{\sigma}
\newcommand{\tc}{\, |\, }

\newcommand{\scalar}[2]{\langle #1, #2\rangle}



\DeclareMathOperator*{\union}{\bigcup}
\DeclareMathOperator*{\inter}{\bigcap}
\DeclareMathOperator{\var}{Var}
\newcommand{\tv}{\mathrm{TV}}
\newcommand{\ent}{\mathrm{Ent}}
\newcommand{\gap}{\mathrm{gap}}
\newcommand{\rBe}{\mathrm{Be}}
\newcommand{\IND}{\mathbf{1}}
\newcommand{\ind}{\mathbf{1}}

\newcommand{\cA}{\mathcal{A}}
\newcommand{\cB}{\mathcal{B}}
\newcommand{\cE}{\mathcal{E}}
\newcommand{\cL}{\mathcal{L}}
\newcommand{\cP}{\mathcal{P}}
\newcommand{\cS}{\mathcal{S}}
\newcommand{\cT}{\mathcal{T}}
\newcommand{\cU}{\mathcal{U}}
\newcommand{\cV}{\mathcal{V}}
\newcommand{\bbC}{\mathbb{C}}
\newcommand{\bbE}{\mathbb{E}}

\newcommand{\bbN}{\mathbb{N}}
\newcommand{\bbP}{\mathbb{P}}
\newcommand{\bbR}{\mathbb{R}}
\newcommand{\bbS}{\mathbb{S}}


%\DeclareMathOperator{\cov}{Cov}
%\DeclareMathOperator{\Hess}{Hess}
%\DeclareMathOperator{\dist}{dist}
%\DeclareMathOperator{\osc}{osc}
%\DeclareMathOperator{\Id}{Id}
%\DeclareMathOperator{\diag}{diag}
%\DeclareMathOperator{\supp}{supp}
%\DeclareMathOperator{\PI}{PI}
%\DeclareMathOperator{\Ent}{Ent}




\let\a=\alpha
\let\b=\beta
\let\d=\delta
\let\e=\varepsilon
\let\g=\gamma
\let\l=\lambda
\let\o=\omega
\let\r=\rho
\let\t=\tau
\let\z=\zeta
\let\G=\Gamma
\let\O=\Omega





%%%%%%%%%%%%%%%%%%%%%%%%%%%%%%%%%%%%%%%%%%%%%%%%%%%%%%%%%%%%%%%%%%%%%%%%%%%%%%%%%%%%%%%%%%%%%%%%%%%
\graphicspath{{./figures/}}

\newcommand*{\mk}{\mkern -1mu}
\newcommand*{\Mk}{\mkern -2mu}
\newcommand*{\mK}{\mkern 1mu}
\newcommand*{\MK}{\mkern 2mu}

\hypersetup{urlcolor=purple, linkcolor=blue, citecolor=red}

\newcommand*{\romanenumi}{\renewcommand*{\theenumi}{\roman{enumi}}}
\newcommand*{\Romanenumi}{\renewcommand*{\theenumi}{\Roman{enumi}}}
\newcommand*{\alphenumi}{\renewcommand*{\theenumi}{\alph{enumi}}}
\newcommand*{\Alphenumi}{\renewcommand*{\theenumi}{\Alph{enumi}}}
\let\oldtilde\tilde
\renewcommand*{\tilde}[1]{\mathchoice{\widetilde{#1}}{\widetilde{#1}}{\oldtilde{#1}}{\oldtilde{#1}}}
\let\oldhat\hat
\renewcommand*{\hat}[1]{\mathchoice{\widehat{#1}}{\widehat{#1}}{\oldhat{#1}}{\oldhat{#1}}}
\let\oldforall\forall
\renewcommand*{\forall}{\mathrel{\oldforall}}
\title{Quantitative multivariate local central limit theorem for irreducible measures on the discrete cube}
\author{Pietro Caputo, Daniel Parisi}
\date{doi:10.5802/ahl.219:published-2024; CC BY 4.0}
\begin{document}
\maketitle
\noindent\textit{Editorial note.} The counted statement is Proposition prop:LCLT of Section 2; its proof is given by the authors in the appendix under the heading at 2193, and the appendix's opening paragraph 2195 is outside the spans because it is a reading guide rather than part of the argument.\par\medskip
\noindent\textit{Editorial note.} The two supporting items (the characteristic-function lemma at 546-552 and Lemma BerEs at 2198-2212) are carried with their complete proofs because the main proof invokes them at 2308 and 2315; no separate count is assigned to either.\par\medskip
\noindent\textit{Editorial note.} The companion theorem 591-604 is carried only as the technique pointer cited at 2290; it is counted as its own card and no part of its proof is needed here.\par\medskip
\noindent\textit{Editorial note.} The source line 2290 begins with a stray typewriter character before the words Set z\_N; this is a native source slip and has been left exactly as in the fixed publisher file, with no mathematical change.\par\medskip
\noindent\textit{Editorial note.} All rates and constants are taken from the fixed publisher source unchanged.\par\medskip
\noindent\textit{Allowed standard prerequisites (editorial).}\begin{itemize}
\item Irreducible and nondegenerate measures on the discrete cube. mu is irreducible if every coordinate xi\_\{i,x\} can be flipped with positive probability while the others are held fixed; irreducibility implies det V\_1 != 0, and the converse fails.. Irreducibility is the hypothesis that makes the characteristic function decay exponentially; nondegeneracy alone does not suffice, as the parity example at 541 and 544 shows.
\item Exponential characteristic-function bound. If mu is irreducible there is c = c(mu) > 0 with |psi(t)| <= exp(-c<t,t>) for all t in [-pi,pi]\textasciicircum{}K.. Bounds the second term of the Fourier split on Q\_\{N,K\} minus A\_N in the main proof.
\item Gaussian covariance normalisation and lattice scaling. V\_N = N V\_1, B\_N = sqrt(det V\_N) = N\textasciicircum{}\{K/2\} sqrt(det V\_1), and L\_N = N\textasciicircum{}\{-1/2\} mu[|V\_1\textasciicircum{}\{-1/2\} xi-bar|\textasciicircum{}3].. Supplies the prefactor and the truncation radius 1/(4L\_N) of the local estimate.
\end{itemize}\medskip
\begin{align}\label{induce}
\si\in\O \mapsto \xi_{i,x} = \ind(\si_i=x)\,,\qquad i=1,\,\dots,\,n\,;\; x\in \{1,\,\dots,\,q_i\}.
\end{align}
That is, $\mu$ is the push forward of $p$ by the above map. Note that we did not include the indicator variable $ \ind(\si_i=0)$ since this is uniquely determined as the indicator of the event $\xi_{i,x}=0$ for all $x\in \{1,\,\dots,\,q_i\}$. When $q_i=1$ for all $i$, then $\O=\{0,1\}^n$ can be identified with $X$, $\si$ with $\xi$, and $\mu$ with $p$.


\subsection{Central limit theorem}

The next results are concerned with the behavior of the sum of independent copies $\xi(1),\,\dots,\,\xi(N)$ of a random variable $\xi$ with values in $X$ and distribution $\mu\in\cP(X)$:
\[
S_N = \sum_{j=1}^N\xi(j).
\]
Thus, $S_N$ is a random vector in $\{0,\,\dots,\,N\}^K$. We use the notation $\scalar{t}{s}=\sum_{i,x}t_{i,x}s_{i,x}$ if $t$ and $s$ are indexed by $i=1,\,\dots,\,n$ and $x=1,\,\dots,\,q_i$. We call $V_1$ the covariance matrix of $\mu$,
\[
V_1(i,x;i',x')=\mu(\xi_{i,x}\xi_{i',x'})-\mu(\xi_{i,x})\mu(\xi_{i',x'}),\qquad i=1,\dots,n\,,\; x=1,\,\dots,\,q_i
\]
Thus $V_1$ is a symmetric nonnegative definite $K\times K$ matrix. If $\det V_1\neq 0$ we say that $\mu$ is \emph{nondegenerate}. The central limit theorem asserts that if $\mu $ is nondegenerate, then as $N\to\infty$ one has the weak convergence
\begin{align}\label{eq:CLT}
\frac1{\sqrt N} \;V_1^{-1/2}\left(S_N - \mu^{\otimes N}(S_N)\right)\;
\longrightarrow\; N(0,\ind_K)\,,
\end{align}
where $\mu^{\otimes N}(S_N)\in[0,N]^K$ is the mean of the vector $S_N$ under the product measure $\mu^{\otimes N}$, and $\ind_K$ denotes the $K\times K$ identity matrix, so that $N(0,\ind_K)$ is the standard normal in $K$ dimensions. Note that when $\mu$ is induced by a measure $p\in\cP(\O)$ as described in~\eqref{induce}, then $\mu^{\otimes N}(S_N)_{i,x}=N\pi_{i,x}$ for all $i,x$, where $\pi_{i,x}$ are the marginals of $p$.

The statement~\eqref{eq:CLT} clearly requires that $\mu$ is nondegenerate. However, one can obtain similar statements in the case of degenerate measures, provided one reduces to the nondegenerate modes by eliminating the degenerate ones. More precisely, one can take the eigenvectors of $V_1$ with nonzero eigenvalues as the new variables. A simple example is obtained if e.g.\ $q_i\equiv 1$ and $\mu$ gives probability $1/2$ to all $1$'s and probability $1/2$ to all $0$'s. Here one simply removes all variables but one.

\subsection{Local central limit theorem}

We will need a local version of the central limit theorem. For this we assume the following stronger notion of nondegeneracy, which we refer to as irreducibility.

\begin{defi}\label{def:irr}
A measure $\mu\in\cP(X)$ is called \emph{irreducible} if for all $i=1,\,\dots, \,n$, for all $x\in\{0,\,\dots,\,q_i\}$, there exists $\xi\in X$ such that $\mu(\xi)$ and $\mu(\xi^{(i,x)})$ are both positive, where $\xi^{(i,x)}$ denotes the vector $\xi$ with the $(i,x)^{\rm th}$ coordinate flipped, that is $\xi^{(i,x)}_{j,y}=\xi_{j,y}$ for all $(j,y)\neq (i,x)$, and $\xi^{(i,x)}_{i,x}=1-\xi_{i,x}$.
\end{defi}

\begin{lemm}\label{lem:nondeg}
If $\mu$ is irreducible then it is nondegenerate. The converse does not hold.
\end{lemm}

\begin{lemm}\label{charfunc}
Suppose $\mu$ is irreducible. Then there exists a constant $c=c(\mu)>0$ such that the characteristic function $\psi(t)=\mu(e^{i\scalar{t}{\xi}})$, $t\in\bbR^K$, satisfies
\begin{align*}
|\psi(t)|\leq e^{-c
\scalar{t}{t} }\,,\qquad \text{for all\;\;} t\in[-\pi,\pi]^K.
\end{align*}
\end{lemm}

\begin{theo}\label{th:eqens}
Suppose $p\in\cP(\O)$ is irreducible and let $\r_N$ be an admissible sequence such that
\begin{align}\label{hypo10}
\scalar{\r_N-\pi}{\r_N-\pi} =O(1/N).
\end{align}
Then for all $k=1,\,\dots,\,N$,
\begin{align}\label{eq:chaos00}
\left\|P_k\g(p,\r_N)-p^{\otimes k}\right\|_\tv \leq \frac{C\,k}{\sqrt N},
\end{align}
for some constant $C=C(p)$. Moreover, when $\r_N=\rhoN$, the canonical tensor product $\g_N(p)$ satisfies the stronger estimate
\begin{align}\label{eq:chaos000}
\left\|P_k\g_N(p)-p^{\otimes k}\right\|_\tv \leq \frac{C\,k}{ N}.
\end{align}
\end{theo}

\begin{prop}\label{prop:LCLT}
Suppose $\mu\in\cP(X)$ is irreducible. Then there exists a finite constant $C=C(\mu)$ such that for all $N\in\bbN$,
\begin{align}\label{eq:LCLT}
\max_{M_N}\;
\left |\;\mu^{\otimes N}\left(S_N=M_N\right) \;-\; \frac{e^{-\frac12\,\scalar{z_N}{z_N}}}{(2\pi N) ^{K/2}
\sqrt {\det V_1} }\;
\right| \leq \frac{C}{ N^{(K+1)/2}},
\end{align}
where
\[
z_N:=\frac1{\sqrt N} \;V_1^{-1/2}\left(M_N - \mu^{\otimes N}(S_N)\right),
\]
and the maximum is over all possible values $M_N\in\{0,\,\dots,\,N\}^K$.
\end{prop}

\begin{proof}
We write
\begin{align*}
|\psi(t)|^2 = \left|\mu\left[e^{i\scalar{t}{\xi}}\right]\right|^2 &=\mu\left[\cos\scalar{t}{\xi}\right]^2 + \mu\left[\sin\scalar{t}{\xi}\right]^2 =\sum_{\xi,\,\xi'\,\in\,X}\mu(\xi)\mu(\xi')\cos\scalar{t}{\xi-\xi'},
\end{align*}
where the last equation uses the identity $\cos(\alpha-\beta) = \sin(\alpha)\sin(\beta) + \cos(\alpha)\cos(\beta).$ If $|\theta|\leq \pi$, then $\cos(\theta)\leq 1-2\theta^2/\pi^2$, and therefore,
\begin{equation*}
\cos\scalar{t}{\xi-\xi'} \leq
\begin{cases}
1-\frac{2\scalar{t}{\xi-\xi'}^2}{\pi^2} & \text{if} \ |\xi-\xi'|_1 = 1 \\
1 & \text{if} \ |\xi-\xi'|_1 \neq 1
\end{cases}
\end{equation*}
where $|\xi-\xi'|_1=\sum_{i,x}|\xi_{i,x}-\xi_{i,x}'|$. Since $|\xi-\xi'|_1=1$ iff $\xi'=\xi^{(i,x)}$ for some $i,x$,
\begin{align*}
|\psi(t)|^2 &\leq 1-\frac{2}{\pi^2}\sum_{i,x}\;\sum_{\xi\,\in\,X}\mu(\xi)\mu\left(\xi^{(i,x)}\right)\,t_{i,x}^2
\leq 1-2c\scalar{t}{t},
\end{align*}
where
\[
c:=\frac1{\pi^2} \;\inf_{i,x}\sum_{\xi\,\in\,X}\mu(\xi)\mu\left(\xi^{(i,x)}\right)
.
\]
The irreducibility of $\mu$ is equivalent to $c>0$. Using $x\leq e^{\frac{1}{2}(x^2-1)}$, $x \in [0,1]$, with $x=|\psi(t)|$, we conclude
\begin{equation*}
|\psi(t)|\leq e^{-c\,\scalar{t}{t}}.
\qedhere
\end{equation*}
\end{proof}

\begin{lemm}\label{BerEs}
Let $\mu\in\cP(X)$ be nondegenerate. For any $t\in\bbR^K$, define
\[
\phi_N(t):=\mu\left(e^{i\left\langle V_N^{-1/2}t,\,\bar{\xi}\right\rangle}\right)^N,
\]
where
\[
\bar{\xi}_{i,x}:=\xi_{i,x}-\mu\left[\xi_{i,x}\right], V_N=N V_1,\quad
\text{and}\quad L_N:=\frac1{\sqrt N}\,\mu\left[\big|V_1^{-1/2}\bar{\xi}\big|^3\right].
\]
Then
\begin{align}\label{beres1}
\left|\phi_N(t)-e^{-\frac{1}{2}\scalar{t}{t}}\right|\leq16L_N|t|^3e^{-\frac{1}3\scalar{t}{t}} \,,\qquad |t|\leq \frac{1}{4L_N}.
\end{align}
\end{lemm}

\begin{proof}
We split the proof into two. First, let us suppose that $ \frac{1}{4L_N}\geq|t|\geq\frac{1}{2}L_N^{-\frac{1}{3}}.$ If one has
\begin{align}\label{beres2}
|\phi_N(t)|^2\leq e^{-\frac{2}{3}\scalar{t}{t}}
\end{align}
then
\begin{align}
\left|\phi_N(t)-e^{-\frac12\scalar{t}{t}}\right|\leq|\phi_N(t)|+e^{-\frac12\scalar{t}{t}}\leq 2e^{-\frac13\scalar{t}{t}}\leq16L_N|t|^3e^{-\frac13\scalar{t}{t}}.
\end{align}
We now show~\eqref{beres2}. Let $\phi(t) := \mu\left[e^{i\scalar t{\bar{\xi}}}\right].$ Define $\Psi:=\Psi_1 - \Psi_2,$ where $\Psi_1$ and $\Psi_2$ are independent and identically distributed as $\bar{\xi}$, so that the characteristic function of $\Psi$ is $|\phi(t)|^2,$ and its covariance matrix is $2V_1.$ Next we use

\begin{align}\label{exprest}
R(x)=e^{ix}-\left(1+ix-\frac{x^2}{2}\right),\qquad |R(x)|\leq \min\left\{|x|^2,\frac{|x|^3}{6}\right\},
\end{align}
for all $x \in \bbR.$ By substituting $x=\scalar{t}{\Psi}$ in~\eqref{exprest} and taking the expectation,
\begin{multline}\label{charfunc1}
|\phi(t)|^2 = 1-\frac{1}{2}\mu\left[\left(\scalar{t}{\Psi}\right)^2\right] + \mu\left[R\left(\scalar{t}{\Psi}\right)\right]
\leq 1-\scalar{t}{V_1t} + \frac{1}{6}\mu\left[\left|\scalar{t}{\Psi}\right|^3\right].
\end{multline}
Then, by using the fact that $(1+x)^N\leq e^{Nx}$ for all $x,$ that $V_N^{-1/2}$ is self-adjoint, and the Cauchy--Schwarz inequality,

\begin{align}
\begin{split}
|\phi_N(t)|^2 =\left|\phi\left(V_N^{-1/2}t\right)\right|^{2N}
&\leq\left(1-\frac1N \scalar{t}{t} + \frac16 \mu\left[\left|\left\langle t,V_N^{-1/2}\Psi\right\rangle\right|^3\right]\right)^N\\
&\leq \left(1-\frac1N \scalar{t}{t} + \frac{8}{6}|t|^3\mu\left[\left|V_N^{-1/2}\bar{\xi}\right|^3\right]\right)^N\\
&\leq \exp\left\{-\scalar{t}{t}+\frac{8}{6}L_N|t|^3\right\} \leq e^{-\frac{2}{3}\scalar{t}{t}},
\end{split}
\end{align}
thus~\eqref{beres1} is proved if $ \frac{1}{4L_N}\geq|t|\geq\frac{1}{2}L_N^{-\frac{1}{3}}.$

Next, suppose that $|t|< \frac{1}{2}L_N^{-\frac{1}{3}}.$ Define $\tau_N := \frac{L_N}{N}.$ One has
\begin{align}\label{ineq}
\frac{1}{2}>L_N^{\frac{1}{3}}|t|>\tau_N^{\frac{1}{3}}|t|=\mu\left[\left|V_1^{-1/2}\bar{\xi}\right|^3\right]^{\frac{1}{3}}\frac{|t|}{\sqrt{N}}\geq \mu\left[\left|V_1^{-1/2}\bar{\xi}\right|^2\right]^{\frac{1}{2}}\frac{|t|}{\sqrt{N}}\geq \frac{|t|}{\sqrt{N}},
\end{align}
where the last inequality follows from
\begin{align}
\mu\left[\left|V_1^{-1/2}\bar{\xi}\right|^2\right]^{\frac{1}{2}} &= \left(\mu\left[\left\langle V_1^{-1/2}\bar{\xi}, V_1^{-1/2}\bar{\xi}\right\rangle \right]\right)^{\frac{1}{2}} = \sqrt{K}\geq 1.
\end{align}
From~\eqref{exprest} and the inequalities~\eqref{ineq}, using again the Cauchy--Schwarz inequality we can write
\begin{multline}\label{beres3}
\left|\phi\left(V_N^{-1/2}t\right)-1\right|
\\
= \left|-\frac1{2N}\scalar{t}{t} + \mu\left[R\left(\left\langle t,V_N^{-1/2}\bar{\xi}\right\rangle\right)\right]\right| \leq \frac{1}{2N}\scalar{t}{t} + \tau_N\frac{|t|^3}{6} < \frac{1}{8}+\frac{1}{48}<\frac{1}{4}.
\end{multline}
For all $|z|<1$ the following inequality holds:
\begin{align}\label{log}
\left|\log(z+1)-z\right| \leq \frac{|z|^2}{2(1-|z|)}
\end{align}
Then, by using, in order, \eqref{log}, \eqref{beres3}, $|a+b|^2\leq2(|a|^2+|b|^2)$ and~\eqref{ineq},
\begin{multline}
\left|\log\phi\left(V_N^{-1/2}t\right)-\left(\phi\left(V_N^{-1/2}t\right)-1\right)\right|\\
\begin{aligned}
&\leq \frac{\left|\phi\left(V_N^{-1/2}t\right)-1\right|^2}{2\left(1-\left|\phi\left(V_N^{-1/2}t\right)-1\right|\right)} <\frac{2}{3}\left|\phi\left(V_N^{-1/2}t\right)-1\right|^2 \\
&=\frac{2}{3}\left|-\frac{1}{2N}\scalar{t}{t} +\mu\left[R\left(\left\langle t,V_N^{-1/2}\bar{\xi}\right\rangle \right)\right] \right|^2 \leq \frac{4}{3}\left(\frac{|t|^4}{4N^2}+ \frac{\tau_{N}^2}{36}|t|^6\right) \\
&\leq \frac{4}{3}\left(\frac{\tau_N}{\sqrt{N}}\frac{|t|^4}{4} + \frac{\tau_{N}^2}{36}|t|^6\right) \leq \frac{4}{3}\left(\frac{1}{2\cdot4}+\frac{1}{36\cdot8}\right)\tau_N|t|^3 =\frac{37}{216}\tau_N|t|^3 < \frac{1}{5}\tau_N|t|^3.
\end{aligned}
\end{multline}
By the triangular inequality we have
\begin{multline}
\left|\log\phi\left(V_N^{-1/2}t\right)+\frac{1}{2N}\scalar{t}{t}\right| \leq \left|\log\phi\left(V_N^{-1/2}t\right) - \left(\phi\left(V_N^{-1/2}t\right)-1\right)\right|\\
+\left|\left(\phi\left(V_N^{-1/2}t\right)-1\right)+\frac{1}{2N}\scalar{t}{t}\right| \leq \frac{1}{5}\tau_N|t|^3 + \frac{1}{6}\tau_N|t|^3 \leq \frac{1}{2}\tau_N|t|^3.
\end{multline}
Using the inequality $|e^z-1|\leq |z|e^{|z|}$, $z \in \bbC$, and~\eqref{ineq},

\begin{multline}
\left|\phi_N(t)e^{\frac{|t|^2}{2}}-1\right| \leq\left|\log \phi_N(t)e^{\frac{|t|^2}{2}}\right|\exp\left(\left|\log \phi_N(t)e^{\frac{|t|^2}{2}}\right|\right) \\
=N\left|\log \phi\left(V_N^{-1/2}t\right)+\frac{1}{2N}\scalar{t}{t}\right|\exp\left(N\left|\log \phi\left(V_N^{-1/2}t\right)+\frac{1}{2N}\scalar{t}{t}\right|
\right) \\
\leq \frac{1}{2}L_N|t|^3e^{\frac{1}{2}L_N|t|^3}
\leq \frac{1}{2}L_N|t|^3e^{\frac{1}{16}} <|t|^3L_N,
\end{multline}
so that~\eqref{beres1} is proved.
\end{proof}

\begin{proof}[{\bf Proof of Proposition~\ref{prop:LCLT}}]
xSet $z_N:=\frac1{\sqrt N} \;V_1^{-1/2}(M_N - \mu^{\otimes N}(S_N))$. Using the Fourier transform as in the proof of Theorem~\ref{th:eqens} and the identity
\begin{align}
\int_{\bbR^K}e^{-i\scalar{s}{z_N}-\frac{\scalar{s}{s}}{2}}ds = (2\pi)^{\frac{K}{2}}e^{-\frac{\scalar{z_N}{z_N}}{2}},
\end{align}
one has
\begin{multline}\label{summation}
\left |\mu^{\otimes N}\left(S_N = M_N\right)-\frac{e^{-\frac{1}{2}\scalar{z_N}{z_N}}}{(2\pi N)^{\frac{K}{2}}\sqrt{\det V_1}}\right| \\
\begin{aligned}
& = \frac{1}{B_N(2\pi)^K}\left | \int_{Q_{N,K}}e^{-i\scalar{s}{z_N}}\phi_N(s)ds-\int_{\bbR^K}e^{-i\scalar{s}{z_N}-\frac{1}{2}\scalar{s}{s}}ds\right| \\
& \leq \frac{1}{B_N(2\pi)^K}\left(\int_{A_N}\left|\phi_N(s)-e^{-\frac{1}{2}\scalar{s}{s}}\right| ds + \int_{Q_{N,K} \setminus A_N}\left|\phi_N(s)\right|ds + \int_{A_N^c}e^{-\frac{1}{2}\scalar{s}{s}}ds\right),
\end{aligned}
\end{multline}
where $B_N:=\sqrt{\det V_N}=N^{K/2}\sqrt {\det V_1}$, $Q_{N,K}:=V_N^{1/2}[-\pi,\pi]^K$, and $A_{N} := \{s\in\bbR^K:\, |s|\leq \frac{1}{4L_N}\}.$ We are going to show that the three terms in the parenthesis above are bounded by $C/\sqrt N$.

Let us define
\[
\tau:=\mu\left[\left|V_1^{-1/2}\bar{\xi}\right|^3\right].
\]
For the first term in~\eqref{summation} we use Lemma~\ref{BerEs}:
\begin{align}
\begin{split}
\int_{A_{N}}\left|\phi_N(s)-e^{-\frac{\langle s,s\rangle}{2}}\right| ds &\leq 16L_N \int_{A_{N}}|s|^3 e^{-\frac{\langle s,s\rangle}{3}}\\
&\leq16L_N \int_{\bbR^K}|s|^3 e^{-\frac{\langle s,s\rangle}{3}} = \frac{C_1}{\sqrt{N}},
\end{split}
\end{align}
where $C_1 :=16\,\tau\int_{\bbR^K}|s|^3 e^{-\frac{<s,s>}{3}}.$ For the second term in~\eqref{summation} we use Lemma~\ref{charfunc}:
\begin{align}\label{eq:t2eq}
\int_{Q_{N,K} \setminus A_{N}}\left|\phi_N(s)\right|ds &= B_N\int_{[-\pi,\pi)^K \setminus \cB_{\tau}}|\phi(t)|^N dt
\leq B_N(2\pi)^Ke^{-\frac{CN}{16\tau^2}},
\end{align}
where $C$ is defined as in that lemma and $\cB_{\tau}:= \{t\in\bbR^K:\,|t|\leq\frac{1}{4\tau}\}$. Therefore, \eqref{eq:t2eq} is bounded by $C_2/\sqrt N$ for a suitable constant $C_2$. Finally, to bound the last term in~\eqref{summation}, using $L_N^{-1}\geq c\sqrt N$ for some constant $c>0$, a simple estimate on the gaussian integral shows that
\begin{align}
\int_{|s|>\frac{1}{4L_N}}e^{-\frac{<s,s>}{2}}ds & \leq \frac{C_3}{\sqrt N},
\end{align}
for some constant $C_3$.
\end{proof}

\begin{thebibliography}{99}
\bibitem{Bertini_etal} Basile, G.; Benedetto, D.; Caglioti, E.; Bertini, L. Large deviations for a binary collision model: energy evaporation, Math. Eng., Volume 5 (2023) no. 1, 1, p. 12
\bibitem{barthe_et_al} Barthe, F.; Cordero-Erausquin, D.; Ledoux, M.; Maurey, B. Correlation and Brascamp–Lieb inequalities for Markov semigroups, Int. Math. Res. Not., Volume 2011 (2011) no. 10, pp. 2177-2216
\bibitem{bonetto2017entropy} Bonetto, F.; Geisinger, A.; Loss, M.; Ried, T. Entropy decay for the Kac evolution, Commun. Math. Phys., Volume 363 (2018) no. 3, pp. 847-875
\bibitem{BT} Bobkov, S. G.; Tetali, P. Modified logarithmic Sobolev inequalities in discrete settings, J. Theor. Probab., Volume 19 (2006) no. 2, pp. 289-336
\bibitem{asymm} Caputo, P.; Martinelli, F. Asymmetric diffusion and the energy gap above the 111 ground state of the quantum XXZ model, Commun. Math. Phys., Volume 226 (2002) no. 2, pp. 323-375
\bibitem{entprod} Caputo, P.; Sinclair, A. Entropy production in nonlinear recombination models, Bernoulli, Volume 24 (2018) no. 4B, pp. 3246-3282
\bibitem{carlenwild} Carlen, E. A.; Carvalho, M. C.; Gabetta, E. Central limit theorem for Maxwellian molecules and truncation of the Wild expansion, Commun. Pure Appl. Math., Volume 53 (2000) no. 3, pp. 370-397
\bibitem{entfound} Carlen, E. A.; Carvalho, M. C.; Le Roux, J.; Loss, M.; Villani, C. Entropy and chaos in the Kac model, Kinet. Relat. Models, Volume 3 (2010) no. 1, pp. 85-122
\bibitem{revchaos} Chaintron, L.-P.; Diez, A. Propagation of chaos: a review of models, methods and applications. I: Models and methods, Kinet. Relat. Models, Volume 15 (2022) no. 6, pp. 895-1015
\bibitem{cattiaux} Cattiaux, P.; Guillin, A.; Malrieu, F. Probabilistic approach for granular media equations in the non-uniformly convex case, Probab. Theory Relat. Fields, Volume 140 (2008) no. 1-2, pp. 19-40
\bibitem{Chung} Chung, K. L. A course in probability theory, Academic Press Inc., 2000
\bibitem{Virag} Csóka, E.; Harangi, V.; Virág, B. Entropy and expansion, Ann. Inst. Henri Poincaré, Probab. Stat., Volume 56 (2020) no. 4, pp. 2428-2444
\bibitem{elementary} Durmus, A.; Eberle, A.; Guillin, A.; Zimmer, R. An elementary approach to uniform in time propagation of chaos, Proc. Am. Math. Soc., Volume 148 (2020) no. 12, pp. 5387-5398
\bibitem{dolera_etal} Dolera, E.; Gabetta, E.; Regazzini, E. Reaching the best possible rate of convergence to equilibrium for solutions of Kac’s equation via central limit theorem, Ann. Appl. Probab., Volume 19 (2009) no. 1, pp. 186-209
\bibitem{DelMoralGuionnet} Del Moral, P.; Guionnet, A. On the stability of interacting processes with applications to filtering and genetic algorithms, Ann. Inst. Henri Poincaré, Probab. Stat., Volume 37 (2001) no. 2, pp. 155-194
\bibitem{miclo} Del Moral, P.; Miclo, L. Genealogies and increasing propagation of chaos for Feynman–Kac and genetic models, Ann. Appl. Probab., Volume 11 (2001) no. 4, pp. 1166-1198
\bibitem{DiacShash} Diaconis, P.; Shahshahani, M. Generating a random permutation with random transpositions, Z. Wahrscheinlichkeitstheor. Verw. Geb., Volume 57 (1981), pp. 159-179
\bibitem{DS} Diaconis, P.; Saloff-Coste, L. Logarithmic Sobolev inequalities for finite Markov chains, Ann. Appl. Probab., Volume 6 (1996) no. 3, pp. 695-750
\bibitem{erbar} Erbar, M.; Fathi, M.; Schlichting, A. Entropic curvature and convergence to equilibrium for mean-field dynamics on discrete spaces, ALEA, Lat. Am. J. Probab. Math. Stat., Volume 17 (2020) no. 1, pp. 445-471
\bibitem{einav2011villani} Einav, A. On Villani’s conjecture concerning entropy production for the Kac master equation, Kinet. Relat. Models, Volume 4 (2011) no. 2, pp. 479-497
\bibitem{Gei} Geiringer, H. On the probability theory of linkage in Mendelian heredity, Ann. Math. Stat., Volume 15 (1944), pp. 25-57
\bibitem{Goel} Goel, S. Modified logarithmic Sobolev inequalities for some models of random walk, Stochastic Processes Appl., Volume 114 (2004) no. 1, pp. 51-79
\bibitem{GQ} Gao, F.; Quastel, J. Exponential decay of entropy in the random transposition and Bernoulli–Laplace models, Ann. Appl. Probab., Volume 13 (2003) no. 4, pp. 1591-1600
\bibitem{grunbaumchaos} Grünbaum, F. A. Propagation of chaos for the Boltzmann equation, Arch. Ration. Mech. Anal., Volume 42 (1971), pp. 323-345
\bibitem{onkacschaos} Hauray, M.; Mischler, S. On Kac’s chaos and related problems, J. Funct. Anal., Volume 266 (2014) no. 10, pp. 6055-6157
\bibitem{KAC} Kac, M. Foundations of kinetic theory, Proceedings of the Third Berkeley Symposium on Mathematical Statistics and Probability, 1954–1955, vol. III, University of California Press (1956), pp. 171-197
\bibitem{lacker} Lacker, D. Hierarchies, entropy, and quantitative propagation of chaos for mean field diffusions, Probab. Math. Phys., Volume 4 (2023) no. 2, pp. 377-432
\bibitem{Lyubich} Lyubich, Y. I. Mathematical structures in population genetics, Biomathematics (Berlin), 22, Springer, 1992 (translated from the Russian by D. Vulis and A. Karpov. Edited by E. Akin)
\bibitem{Martinez} Martínez, S. A probabilistic analysis of a discrete-time evolution in recombination, Adv. Appl. Math., Volume 91 (2017), pp. 115-136
\bibitem{mckean1966speed} McKean, H. P. jun. Speed of approach to equilibrium for Kac’s caricature of a Maxwellian gas, Arch. Ration. Mech. Anal., Volume 21 (1966), pp. 343-367
\bibitem{mckean} McKean, H. P. jun. An exponential formula for solving Boltzmann’s equation for a Maxwellian gas, Comb. Theory, Volume 2 (1967), pp. 358-382
\bibitem{kacsprog} Mischler, S.; Mouhot, C. Kac’s program in kinetic theory, Invent. Math., Volume 193 (2013) no. 1, pp. 1-147
\bibitem{anewapp} Mischler, S.; Mouhot, C.; Wennberg, B. A new approach to quantitative propagation of chaos for drift, diffusion and jump processes, Probab. Theory Relat. Fields, Volume 161 (2015) no. 1-2, pp. 1-59
\bibitem{Nagylaki} Nagylaki, T. Introduction to theoretical population genetics, Biomathematics (Berlin), 21, Springer, 2013
\bibitem{Petrov} Petrov, V. V. Sums of independent random variables, Walter de Gruyter, 2022
\bibitem{review_discrete} Platkowski, T.; Illner, R. Discrete velocity models of the Boltzmann equation: a survey on the mathematical aspects of the theory, SIAM Rev., Volume 30 (1988) no. 2, pp. 213-255
\bibitem{Reza} Rezakhanlou, F. Propagation of chaos for particle systems associated with discrete Boltzmann equatio, Stochastic Processes Appl., Volume 64 (1996) no. 1, pp. 55-72
\bibitem{Sinetal2} Rabani, Y.; Rabinovich, Y.; Sinclair, A. A computational view of population genetics, Random Struct. Algorithms, Volume 12 (1998) no. 4, pp. 313-334
\bibitem{Sinetal} Rabinovich, Y.; Sinclair, A.; Wigderson, A. Quadratic dynamical systems, 33rd annual symposium on Foundations of computer science (FOCS)Proceedings, IEEE (1992), pp. 304-313
\bibitem{Baakeetal} Salamat, M.; Baake, M.; Baake, E. The general recombination equation in continuous time and its solution, Discrete Contin. Dyn. Syst., Volume 36 (2016) no. 1, pp. 63-95
\bibitem{sznitman} Sznitman, A.-S. Topics in propagation of chaos, Calcul des probabilités, Ec. d’Été, Saint-Flour/Fr (Lecture Notes in Mathematics), Volume 1464, Springer (1989), pp. 165-251
\bibitem{tanaka} Tanaka, H. Probabilistic treatment of the Boltzmann equation of Maxwellian molecules, Z. Wahrscheinlichkeitstheor. Verw. Geb., Volume 46 (1978), pp. 67-105
\bibitem{Vil} Villani, C. Cercignani’s conjecture is sometimes true and always almost true, Commun. Math. Phys., Volume 234 (2003) no. 3, pp. 455-490
\bibitem{Wild} Wild, E. On Boltzmann’s equation in the kinetic theory of gases, Proc. Camb. Philos. Soc., Volume 47 (1951) no. 3, pp. 602-609
\end{thebibliography}
\end{document}
